\section{Approach}
\label{sec:approach}

We now describe our implementation of relational e-matching in Sundance, an SMT
solver currently in development~\cite{sundance}. The implementation spans three
components: (1)~two index data structures that map e-class roots to the e-nodes
that reference them, (2)~a pattern-flattening compiler that reduces a nested
trigger to a conjunctive query over these indices, and (3)~a left-deep join
executor with greedy query planning, and semi-naive evaluation.

\subsection{Index data structures}
\label{sec:indices}

The core idea of relational e-matching is to precompute indices so that, given a
bound variable's e-class, one can directly look up which e-nodes reference that
e-class---without scanning every e-node of a given function symbol.

Sundance maintains two such indices, keyed by function symbol.

\paragraph{Output index.} For each function symbol $f$, the \emph{output index}
maps an e-class root $c$ to the set of $f$-e-nodes whose output (i.e., the
e-class of the term $f(\ldots)$ itself) is~$c$. In code, \texttt{function\_maps}
is itself keyed by function symbol, with type
%
\[
  \text{FnSymbol} \;\to\; (\text{e-class} \to \{\text{e-node uid}\}).
\]

\paragraph{Argument index.} For each function symbol $f$ and each argument
position $i$, the \emph{argument index} maps an e-class root $c$ to the set of
$f$-e-nodes whose $i$-th argument is in e-class~$c$. In code,
\texttt{function\_indices} has type
%
\[
  \text{FnSymbol} \;\to\; (\text{ArgPos} \to (\text{e-class} \to \{\text{e-node uid}\}))
\]

Both indices are implemented as \texttt{im::OrdMap} (a persistent balanced tree
from the \texttt{im} crate~\cite{im-rs}). When two e-classes are merged, entries from the old root are moved
to the new root in both indices. Each entry carries a timestamp (the
\emph{matching round} at which it was inserted) to support the semi-naive
evaluation discussed in Section~\ref{sec:seminaive}.

\paragraph{Backtracking.} At each SAT decision level, Sundance saves a snapshot
of the two indices. Because \texttt{im::OrdMap} is a persistent data structure,
cloning is $O(1)$ --- it shares structure with the live copy. When backtracking, the
solver restores the snapshot, discarding all index mutations that occurred under
the popped decision level. New terms created by quantifier instantiations are
then re-inserted into the restored indices, since those terms are
permanent in the e-graph even though the equalities that triggered them may have
been undone.

The cheap snapshots come at a cost: \texttt{im::OrdMap} is a persistent balanced
tree and every insert, lookup, or removal pays the usual $O(\log n)$ tree walk. We 
discuss alternative implementations in Section~\ref{sec:conclusion}.

\subsection{Pattern flattening}
\label{sec:flattening}

A trigger pattern such as \smt{(g (h x) x)} is a nested term. Our relational
matcher requires it to be flattened into a conjunction of relational atoms, such 
as $h(x) = v_0  \wedge  g(v_0,\; x) = root$

The resulting atom list is a conjunctive query: a match must simultaneously
satisfy all atoms. Each quantifier-bound variable (e.g., \smt{x}) may appear in
multiple atoms, acting as a join variable. Ground sub-terms (constants) are
recorded as fixed bindings and used to constrain index lookups.

\subsection{Join execution}
\label{sec:join}

Given the flattened atom list, the join executor produces all variable
assignments that satisfy every atom simultaneously. We use a left-deep binary
join: atoms are processed one at a time in a fixed order, and at each step the
current set of partial bindings (the \emph{frontier}) is extended by joining
against the next atom.

\subsubsection{Query planning}

A good atom ordering avoids Cartesian products by ensuring that each new atom
shares at least one variable with the frontier. Sundance uses a greedy heuristic:
\begin{enumerate}
\item Pick the atom whose function has the fewest e-nodes (smallest table).
\item For each subsequent position, pick the remaining atom that shares the most
      variables with the already-bound set. Break ties by smallest table size.
\end{enumerate}
If any atom's function has zero e-nodes, the join is guaranteed empty and we
short-circuit immediately.

\subsubsection{Candidate retrieval}

At each join step, for each partial binding in the frontier, we must find all
e-nodes of the current atom's function that are consistent with the bound
variables. This is where the index data structures pay off.

For each argument position $i$ of the atom whose variable is already bound to
some e-class~$c$, we query \linebreak $\text{function\_indices}[f].\text{args}[i][c]$
to retrieve candidate e-node uids. If the output variable is bound, we similarly
query the output index. When multiple variables are bound, we intersect the
candidate sets.
%  (via a sorted-merge intersection).
If no variables are bound (as
with the first atom in the join order), we fall back to a full scan of the
function's e-node table.

% Crucially, the index lookups are hoisted out of the per-binding loop: for a
% given join step, the function-name lookups into \texttt{function\_maps} and
% \texttt{function\_indices} are performed once and reused across all frontier
% bindings.

% \subsubsection{Binding extension and consistency}

% For each candidate e-node, we attempt to extend the current partial binding. For
% each variable position (argument or output), if the variable is already bound, we
% check that the candidate's e-class at that position equals the bound value
% (using the union-find's \texttt{find} operation). If the variable is unbound, we
% bind it to the candidate's raw term uid and record the canonical e-class root
% alongside it.

% Storing the raw uid (rather than the canonical root) is important for
% correctness: the downstream quantifier instantiation substitutes these terms
% into the quantifier body, and using the raw fnode argument preserves the
% invariant that two matches against the same fnode produce identical
% substitutions. This prevents the solver from re-instantiating semantically
% equivalent quantifier bodies with different class representatives across rounds.

% \subsubsection{Deduplication}
% \label{sec:dedup}

% Multiple e-nodes in the same e-class can produce bindings that differ in raw
% uids but agree on canonical roots.  Without deduplication, these redundant
% bindings cause a Cartesian blowup: if a function has $m$ e-nodes all in the
% same e-class, the next join step sees $m$ ``copies'' of what is semantically one
% binding, multiplying the frontier by $m$ at every step.

% To avoid this, after each join step Sundance sorts the new bindings by their
% canonical-root vectors and removes duplicates. Because we store the canonical
% root alongside the raw uid at bind time, this deduplication requires no
% additional \texttt{find} calls---it is a simple vector comparison and dedup.

\subsection{Semi-naive evaluation}
\label{sec:seminaive}

Each matching round assigns a timestamp to newly inserted index entries. Between
rounds, the \emph{matching round counter} is incremented. Semi-naive evaluation
ensures that each round only discovers \emph{new} matches by requiring at least
one atom in the join to use only ``delta'' entries (those with timestamp
$\geq$ the current matching round). For a flattened pattern with $k$ atoms, we
execute $k$ passes: in pass $i$, atom $i$ is restricted to delta entries while
all other atoms use the full index. The union of results across all passes yields
exactly the new matches. If no atom has any delta entries, the entire
multipattern is skipped.

% \subsection{Discussion: room for improvement}
% \label{sec:room}

% Our current implementation is a first step; several optimizations remain.

% \paragraph{Pattern flattening does not share subterms.}
% If the same sub-expression appears twice in a trigger (e.g.,
% \smt{f(g(x), g(x))}), the flattener creates two independent fresh variables for
% the two occurrences of \smt{g(x)} rather than reusing one. This can introduce
% redundant atoms and inflate the join.

% \paragraph{String-keyed lookups in the hot path.}
% Function symbols are currently stored as heap-allocated strings. Every index
% lookup hashes the function name; interning symbols to integer IDs would replace
% these with array lookups.

% \paragraph{Per-quantifier setup cost.}
% For each quantifier in each matching round, we rebuild the variable-index map
% and recompute the join order. Pre-computing these at pattern-registration time
% and caching them on the quantifier would eliminate this per-round overhead.

% \paragraph{Post-match pipeline dominance.}
% Our profiling shows that the matching phase itself (index lookups and joins) is
% often a small fraction of the total per-round cost. The bulk of the time is
% spent in the post-match pipeline: substituting the matched bindings into the
% quantifier body, converting to NNF, inserting the resulting terms into the
% egraph (\texttt{insert\_predecessor}), and converting to CNF. These costs are
% shared between relational and top-down e-matching and are orthogonal to the
% matching algorithm, but they currently dominate wall-clock time on realistic
% benchmarks.

% \paragraph{Worst-case optimal joins.}
% Our current join executor uses a left-deep binary join, which is not worst-case
% optimal for cyclic queries. Adopting a worst-case optimal join algorithm (e.g.,
% generic join~\cite{genericjoin}) could improve performance on patterns where the
% optimal plan is not left-deep.
